\section{3D 绘制路径复查}\label{sec:3d-review}
这次复查从同一个物体在不同绘制阶段的行为入手。以下数值来自独立 GPU 测试，原有 RenderDoc 捕获仍只对应深度预传递与普通主绘制，不能代替本节新增证据。

\subsection{屏幕外物体被一并移出阴影}
原先在主相机视锥外就丢弃绘制命令。相机看不到的物体仍可能处在光源范围内，给可见表面投下阴影或贡献 RSM 数据。过早删除命令，会让这些影响一起消失。

现在保留提交内容，主相机的深度预传递与主绘制按相机可见性筛选，Shadow 与 RSM 仍处理该物体。测试把四边形移到相机外，再调整光源矩阵使其落在光源范围内。主画面不写深度，阴影深度和 RSM 表面数据仍出现，绘制次数也与预期一致。另检查相交边界和负缩放包围盒。
\loc{Redflag/src/Redflag/Renderer/Renderer3D.cpp}
\loc{tests/Renderer3DTests.cpp}

\subsection{阴影与 RSM 的透明裁剪缺项}
修复普通主绘制和深度预传递后，阴影与 RSM 仍需逐条核对。它们此前没有使用完整的底色因子、底色贴图 alpha 与透明度图红通道乘积，会留下本该裁掉的遮挡面。

现在三项输入共用同一材质上传规则。Renderer3DTests 对 $32\times32$ 的四条纹逐像素检查 Shadow 深度和 RSM 深度、表面有效标记。Mask 只保留两侧，Opaque 保留全部，Blend 仍跳过这两个光源绘制阶段。测试通过不表示半透明投影已经实现。
\loc{RedflagForge/assets/shaders/ShadowDepth.glsl}
\loc{RedflagForge/assets/shaders/RSMDepth.glsl}

\subsection{普通与实例化主绘制的光照不一致}
两份主 Shader 曾分别维护光照，实例化分支可能遗漏普通路径已经加入的 RSM 输入。现在用同一个源码和实例化编译选项组织两种顶点输入，片元光照共用实现。

测试输入固定红色间接光，普通与实例化输出的红通道都增加 0.5，绿蓝不变。实例化用一次绘制写出左右实体 ID 71、72。这证明指定输入下两条路径读取了相同的间接光规则，没有测量性能提升。当前是 15 个 .glsl 文件，另编译一个实例化版本。
\loc{tests/DepthPrepassTests.cpp}

\subsection{RSMDepth 的跨阶段参数块重名}
新测试第一次链接失败时，顶点和片元阶段各自生成了临时名 \code{\_11\_13}，实际内容却分别是模型矩阵与材质参数。链接器将同名块视为同一声明，产生冲突。这次问题有实际链接失败记录，与旧 varying 名称问题不同。

反编译阶段现在按资源绑定位置统一命名块类型与实例，使同一份数据名称一致、不同数据名称不同。修复后 Shader 链接与 Renderer3D 逐像素测试通过。这个规则服务当前 OpenGL 布局，不是 Vulkan 后端已完成的证明。
\loc{Redflag/src/Platform/OpenGL/OpenGLShader.cpp}

\newpage
\section{导入数据与 IBL 资源复查}
\subsection{显式零值被当成缺失参数}
旧导入逻辑会在有贴图时将接近黑色的底色改白，并把零金属度、零粗糙度视为没有提供参数。零是合法材质值，不能用数值本身判断属性是否存在。

现在根据 Assimp 查询是否成功区分缺失与显式值，保留黑色、底色 alpha 和两个零因子。新增 glTF 用例同时给出白色贴图与 $(0,0,0,0.25)$ 底色，以及金属度、粗糙度零值，导入后逐项断言保留。模型测试因此从九例扩为十例。无对应贴图且属性缺失时仍使用金属度 0、粗糙度 0.5，有贴图而缺少乘数时使用 1。

这验证的是导入数据不被猜测覆盖。材质仍没有完整 Blend 导入与排序，保存 alpha 不能被描述为玻璃已经正确显示。
\loc{tests/ModelImportTests.cpp}

\subsection{IBL 预计算与绘制状态}
预计算和全屏后处理会在场景绘制之间运行。继承剔除、裁剪或拒绝所有片元的深度状态，会使预计算输出为空。现在显式设置必要状态，结束后恢复。IBLTests 故意设置这些状态，再执行真实 HDR 预计算，检查产物和调用后状态。

另一条失败路径在 HDR 读取前就创建辅助四边形，Shutdown 却没有删除它。现在成功解码后才分配，删除无调用方的空环境准备接口，Shutdown 同时释放四边形与立方体资源。

测试使用常量 HDR $(2,1,0.5)$，读取环境图与辐照度图的六面，以及预过滤图五个 mip 层。颜色在容差内保持一致，BRDF 查找表为有限值且有有效输出。预计算后 viewport 与相关绘制状态恢复，Shutdown 后产物纹理和辅助 VAO/VBO 不再有效。缺失文件也保持未就绪并释放资源。

后处理测试另在剔除全部面、只允许 $1\times1$ 裁剪区域的条件下绘制，确认完整目标仍有均匀输出，调用后保留外部开关。原有双实例曝光和尺寸检查继续通过。测试只覆盖所列状态与输入，没有做全部 GL 状态组合或长时间压力测试。
\loc{tests/IBLTests.cpp}
\loc{tests/PostProcessTests.cpp}

\subsection{纹理上传的行对齐}
宽度为 3 的 RGB 图每行有 9 字节，文件解码数据紧密排列，默认四字节对齐却会让下一行从错误位置开始。上传现在临时使用一字节对齐，完成后恢复调用者原值。文件加载和 SetData 使用相同约定，并补齐灰度、灰度加 alpha 的格式与采样分量。

测试对 $3\times2$ 的 1、2、3、4 通道图分别上传、修改并逐像素读回，进入时刻意把对齐设为 8，结束后仍为 8。全部通过。这是新生成素材的结果，不据此给早年棋盘倾斜补写根因。
\loc{Redflag/src/Platform/OpenGL/OpenGLTexture.cpp}

\newpage
\section{场景帧与附件状态复查}
\subsection{深度附件查询与预览}
主目标使用纯深度 DEPTH32F，SSAO 和 RSM 却查询深度模板组合附件。现在统一查询深度附件。深度预览同时修复上传行对齐的恢复，以及整张图只有一种几何深度时显示为空的处理。Shutdown 清掉借用的纹理 ID，避免面板继续使用失效名字。

SceneRenderTests 在编辑和运行路径连续绘制，确认 SSAO/RSM 不产生 GL 错误，历史输出存在，预览中单深度平面可见，销毁后查询返回零。这里验证附件接通和状态，不评价复杂场景的遮蔽效果或历史拖影。

\subsection{跨帧 2D 混合与方向光选择}
后处理关闭混合后，下一帧 2D 没有重新启用，半透明精灵会失去预期混合。场景现在在 2D 阶段设置标准 alpha 混合，在 3D 数据阶段关闭混合，并避免 2D 向法线附件写入未定义值。连续帧测试读到预期的 0.5 颜色和零法线。

方向为零的灯光也曾先归一化再判断，可能产生 NaN。现在先检查向量，再上传有效灯光。主绘制只支持一盏方向光时，阴影必须跟随实际上传的同一盏，不能另找后面的可投影灯。测试分别检查零方向不上传，以及当前灯不投影时阴影关闭。
\loc{tests/SceneRenderTests.cpp}

\subsection{天空盒串写实体号和法线}
这次实测先把实体附件清为 -1、法线首分量清为 0.25，画天空盒后分别变成 1073741824 和 2。天空盒 Shader 只声明颜色输出，但其余附件仍允许写入，不能依赖未声明输出替自己保留旧数据。

修复在天空盒绘制期间关闭实体号和法线附件的写入，再恢复原掩码。最终测试保留 -1 与 0.25，也检查此前颜色写掩码没有改变。它避免天空背景污染拾取与 SSAO 输入，没有增加天空盒参与这些附件的含义。
\loc{Redflag/src/Redflag/Renderer/IBLPrecompute.cpp}
\loc{tests/IBLTests.cpp}

\subsection{复现入口与最终记录}
最终 Debug/x64 构建及八组回归全部通过。新增三组从仓库根目录运行，基础五组仍使用前文脚本。
\begin{verbatim}
tests\run_renderer3d_tests.cmd
tests\run_ibl_tests.cmd
tests\run_scene_tests.cmd
\end{verbatim}
失败与八组通过记录已节选到 \path{output/validation/opengl-regression.txt}，可随仓库复查。本机完整日志为 \path{tmp/comment-clarity/verified-*.log}。15 个 .glsl 独立链接通过，实例化版本另有实际绘制断言。这些专项结果仍不能替代完整编辑器交互、所有资产和长期压力测试。
